\documentclass{beamer}

\usepackage{amsmath}
\usepackage{bm}
\usepackage{tcolorbox}

\title{Topic 1-practice}
\author{Daniel Trad}
\date{}

\begin{document}

%------------------------------------------------
\frame{\titlepage}

%------------------------------------------------
\begin{frame}{Pressure and Displacement Formulations}
\small

Consider acoustic wave propagation in a homogeneous fluid.

\vspace{0.2cm}

The wavefield may be described either by the scalar pressure field
\[
\frac{\partial^2 p}{\partial t^2}=v^2\nabla^2 p,
\]
or by the particle displacement vector $ \bm{u}(\bm{x},t)=\begin{bmatrix}u_x \\ u_y \\ u_z\end{bmatrix}.$

\begin{tcolorbox}[title=Given information,colback=blue!5!white,colframe=blue!60!black]
Assume that pressure is related to displacement through
\[
p=-K\,\nabla\cdot\bm{u},
\]
where \(K\) is the bulk modulus, and $\nabla\cdot\bm{u}=\frac{\partial u_x}{\partial x}+\frac{\partial u_y}{\partial y}+\frac{\partial u_z}{\partial z}.$
\vspace{0.15cm}
Also assume that each displacement component satisfies
\[
\frac{\partial^2 u_i}{\partial t^2}=v^2\nabla^2 u_i,
\qquad i=x,y,z.
\]
\end{tcolorbox}

\end{frame}

%------------------------------------------------
\begin{frame}{Discussion Questions}
\small

Discuss the following:

\begin{enumerate}
    \item Why is \(p\) a \textbf{scalar} field, while \(\bm{u}\) requires \textbf{three components} in 3D?
    
    \vspace{0.15cm}
    
    \item Take the divergence of the three displacement equations. Show that \(\nabla\cdot\bm{u}\) satisfies a scalar wave equation, and therefore so does \(p\).
    
    \vspace{0.15cm}
    
    \item What information about the displacement field is retained by \(\nabla\cdot\bm{u}\), and what information is lost?
    
    \vspace{0.15cm}
    
    \item Why is pressure often a more economical variable for describing acoustic waves in a fluid?
    
    \vspace{0.15cm}
    
    \item Does knowing \(p(\bm{x},t)\) uniquely determine all three components \(u_x,u_y,u_z\)? Explain.
\end{enumerate}

\vspace{0.25cm}

\begin{block}{Conclusion to discuss}
In what sense can one scalar pressure equation replace three displacement equations, and in what sense are the two formulations not completely identical?
\end{block}

\end{frame}

%------------------------------------------------
\begin{frame}{Instructor Solution / Sketch}
\small

\begin{enumerate}
    \item \textbf{Scalar vs. vector:}  
    Pressure \(p\) has one value at each point, so it is scalar.  
    Displacement \(\bm{u}=(u_x,u_y,u_z)\) has three components, so it is vector-valued.

    \vspace{0.15cm}

    \item \textbf{Take the divergence of the displacement equations:}
    \[
    \frac{\partial^2 u_i}{\partial t^2}=v^2\nabla^2 u_i
    \qquad (i=x,y,z).
    \]
    Apply \(\nabla\cdot\):
    \[
    \frac{\partial^2}{\partial t^2}(\nabla\cdot\bm{u})
    =
    v^2\nabla^2(\nabla\cdot\bm{u}).
    \]
    Since
    \[
    p=-K\,\nabla\cdot\bm{u},
    \]
    and \(K\) is constant, it follows that
    \[
    \frac{\partial^2 p}{\partial t^2}=v^2\nabla^2 p.
    \]

    \vspace{0.15cm}

\end{enumerate}
\end{frame}

%------------------------------------------------
\begin{frame}{Instructor Solution / Sketch}
\small

\begin{enumerate}
    \setcounter{enumi}{2}


    \item \textbf{What is retained / lost?}  
    The divergence keeps the \textbf{compressional (volumetric)} part of motion.  
    It loses the full directional information of \(\bm{u}\), especially any part not seen in \(\nabla\cdot\bm{u}\).

    \vspace{0.15cm}

    \item \textbf{Why pressure is economical:}  
    For acoustic waves in fluids, the main physics is compressional, so one scalar \(p\) is often enough instead of three unknowns \(u_x,u_y,u_z\).

    \vspace{0.15cm}

    \item \textbf{Important subtlety:}  
    Knowing \(p\) does \emph{not} generally determine the full displacement vector uniquely.  
    Pressure determines \(\nabla\cdot\bm{u}\), not all components of \(\bm{u}\).

\end{enumerate}

\end{frame}



%------------------------------------------------
\begin{frame}{Reflection at the Free Surface and Multiples}
%\footnotesize
\scriptsize
%\tiny
Consider an upgoing primary wave reaching the Earth's surface.

The pressure reflection coefficient for normal incidence is

\[
R=\frac{Z_2-Z_1}{Z_2+Z_1},
\]

where \(Z=\rho v\) is the acoustic impedance. At the free surface, 

$Z_{\rm air}=\rho_{\rm air}v_{\rm air}\ll Z_{\rm earth},$
and therefore $R_{\rm surface}=\frac{Z_{\rm air}-Z_{\rm earth}}{Z_{\rm air}+Z_{\rm earth}}\approx -1.$

\begin{block}{Physical interpretation}
If the upgoing primary arriving at the surface has amplitude \(P\),

\[
P_{\rm reflected}
=
R_{\rm surface}P
\approx -P.
\]

Thus, the wave reflected back into the Earth has approximately
\textbf{the same amplitude as the arriving primary}, but with
\textbf{opposite polarity}.
\end{block}

This downgoing wave can reflect again from subsurface interfaces and
return to the receivers as a \textbf{free-surface multiple}:

\[
\boxed{
P
\quad\xrightarrow{\text{free surface}}\quad
-P
\quad\xrightarrow{\text{subsurface reflection}}\quad
M
}
\]

Therefore, the free surface does not significantly attenuate the wave:
it sends essentially the full primary amplitude back into the Earth,
allowing strong multiples to be generated.

\end{frame}

\end{document}